\origpage{25}{GRAPH MINORS. XIV}
\hspace*{12pt}(i) \textit{$H \cap (\Sigma - \Lambda)$ is a subdivision of a simple 3-connected graph}

\hspace*{12pt}(ii) \textit{for every $y \in \Lambda$ there is a closed curve $\phi$ meeting $H$ at most 39 times, with $\{x, y\} \subseteq \operatorname{ins}(\phi)$}

\hspace*{12pt}(iii) \textit{$\operatorname{bd}(\Delta) - \Lambda$ meets $H$ at least 4 times, for every closed disc $\Delta \subseteq \Sigma$ with $a, b, c, d \in \Delta$, and}

\hspace*{12pt}(iv) \textit{there are two disjoint paths of $G$ with ends $a, c$ and $b, d$ each meeting $H$ only within $\Lambda \cup \{a, b, c, d\}$.}

This simplified form is not very far from what we really want. The main changes we need are:

\hspace*{12pt}(i) we need to permit $\Sigma$ to be a sphere; and then the definition of $\operatorname{ins}(\phi)$ becomes more complicated, using tangles

\hspace*{12pt}(ii) $G$ might not be 4-connected; then we ask (in (1.1)(i)) only that the drawing be extendible up to 3-separations (this is called the existence of ``a 3-segregation of $G$ and an arrangement of it in $\Sigma$ compatible with $H$'').

The second change actually causes much more trouble then the first. When we make both of these changes, a new technical problem arises---we have to ensure that the 3-separations of $G$ that we use have $H$ on the proper side, that the baby does not get thrown out with the bath-water. To do so, we introduce ``$\mathscr{T}$-local'' 3-segregations. The final version of (1.1) is stated in (4.1), and sections 2 and 3 are concerned with definitions.

The proof we give of our main result is not the simplest possible. We have chosen a proof which involves proving and applying a number of fairly complicated lemmas about ``crooked transactions''; because we shall need these lemmas in any case in the next paper of the series, and, given these lemmas, deriving our present result is easy.

\sectionhead{2. Surfaces}

In this paper, a \textit{surface} is a compact connected 2-manifold, possibly with boundary. In fact, the only surfaces with boundary that we shall need to consider are \textit{closed discs} and \textit{closed cylinders}, homeomorphic to
\begin{equation*}
\{(x, y) \in \mathbf{R}^2 \colon x^2 + y^2 \le 1\}
\end{equation*}
\begin{equation*}
\{(x, y) \in \mathbf{R}^2 \colon 1 \le x^2 + y^2 \le 2\}
\end{equation*}
respectively. \textit{Open discs} are defined similarly (but they are not surfaces). If $\Sigma$ is a surface, its boundary is denoted by $\operatorname{bd}(\Sigma)$; and if $\Lambda \subseteq \Sigma$ is an open disc, its boundary $\operatorname{bd}(\Lambda)$ is defined in the natural way. (We shall use
